%% Proof annex fragment for Part 2.2 calculi.  Rules remain in the registry.

\begin{annexitem}{annex:p2.fol.calc.1}{thm:p02-reduct-invariance}{Invariance under reduct}
\begin{proof}
Term induction gives equal values because every old function symbol has the
same interpretation in the expansion and reduct.  Formula induction then
uses that result for atoms, unchanged old relations, the connective clauses,
and the common domain in each quantifier case.
\end{proof}
\end{annexitem}

\begin{annexitem}{annex:p2.fol.calc.2}{thm:p02-fol-nd-soundness}{Soundness of first-order natural deduction}
\begin{proof}
Induct on the derivation, strengthening the claim to every structure and
assignment satisfying its open assumptions.  Propositional rules use
Theorem~\ref{thm:prop-nd-soundness}.  Universal elimination and existential
introduction follow from the substitution lemma.

For \(\forall I\), update the assignment by arbitrary
\(a\in\dom\mathcal M\).  Eigenvariable freshness and coincidence preserve
all open assumptions; induction gives the matrix, hence the universal.
For \(\exists E\), choose a witness \(a\) and put
\(s'=s[y\mapsto a]\).  Freshness, coincidence, and substitution make the
minor assumption true under \(s'\) and preserve every surviving assumption.
Induction yields its conclusion; since \(y\) is absent there, coincidence
returns to \(s\).

Alpha conversion uses Corollary~\ref{cor:alpha-invariance}.
Identity introduction is reflexivity.  In identity elimination, equal term
values make the two substitution updates identical, so
Theorem~\ref{thm:fol-substitution} transfers satisfaction.  These are all
new rule cases.
\end{proof}
\end{annexitem}

\begin{annexitem}{annex:p2.fol.calc.3}{thm:p02-fol-lk-soundness}{Soundness of the first-order sequent calculi}
\begin{proof}
Induct on the sequent derivation.  Propositional cases use the P2.1 soundness
Theorem~\ref{thm:prop-lk-soundness}; \(\forall L,\exists R\) use term
evaluation plus substitution.

For the repaired \(\forall R\) case, suppose its conclusion is false under
\(\mathcal M,s\).  Thus all antecedent formulas are true, all old succedent
formulas false, and some \(a\in\dom\mathcal M\) falsifies \(\varphi\) under
\(s[x\mapsto a]\).  Put \(s'=s[y\mapsto a]\).  Substitution says that
\(\varphi[y/x]\) under \(s'\) has the value of \(\varphi\) under
\(s'[x\mapsto a]\).  If \(y=x\), this is the chosen update.  Otherwise
eigenvariable freshness gives \(y\notin\FV(\varphi)\), so coincidence again
reduces it to the chosen update.  The same freshness condition puts \(y\)
outside every formula of both contexts; coincidence therefore preserves the
antecedent-true/succedent-false pattern under \(s'\).  The premise is false,
contradicting induction.  The \(\exists L\) argument is dual.

The equality initial sequent is valid by reflexivity; equality replacement
is valid because equal terms induce the same assignment update in the
substitution lemma.  Hence every registered rule preserves sequent validity.
\end{proof}
\end{annexitem}

\begin{annexitem}{annex:p2.fol.calc.4}{thm:p02-equality-laws}{Derived equality laws}
\begin{proof}
Reflexivity and formula replacement are the two registered equality rules.
Choose \(z\notin
\operatorname{Var}(r)\cup\operatorname{Var}(t)\cup\operatorname{Var}(u)\).
For symmetry, take \(r=r\) and apply replacement to \(r=t\) with
\(\varphi(z)\defeq z=r\), obtaining \(t=r\).  For transitivity, apply
replacement to \(t=u\) and \(r=t\) with
\(\varphi(z)\defeq r=z\), obtaining \(r=u\).  Total capture-avoiding
substitution makes both displayed instances legitimate.
\end{proof}
\end{annexitem}

\begin{annexitem}{annex:p2.fol.calc.5}{thm:p02-hilbert-fol-deduction}{Hilbert-FOL deduction theorem}
\begin{proof}
The reverse implication appends \(\varphi\) and applies modus ponens.  For
the forward implication, induct through a dependency-annotated derivation
from \(\Gamma\cup\{\varphi\}\), proving
\(\Gamma\derives_{\mathsf H_{\mathsf{FOL}}}\varphi\to\chi_i\).
Premises from \(\Gamma\) and axioms use A1; the line \(\varphi\) uses the
standard A1/A2 derivation of \(\varphi\to\varphi\); modus ponens uses A2.

Suppose \(\chi_i=\forall x\,\psi\) comes by generalization.  If the
dependency set for \(\psi\) contains \(\varphi\), the rule condition gives
\(x\notin\FV(\varphi)\).  Generalize the induction result
\(\varphi\to\psi\), then use the \(\forall2\) schema to obtain
\(\varphi\to\forall x\psi\).  If \(\varphi\) is absent from that dependency
set, its traced subderivation is already from \(\Gamma\); generalize
\(\psi\) with the same dependencies and use A1 to infer
\(\varphi\to\forall x\psi\).  Thus every annotated line is covered, and the
last one proves the claim without unrestricted generalization.
\end{proof}
\end{annexitem}
